\chapter*{Resumen}
~\newline~\newline~
\markboth{Resumen}{Resument}
\addcontentsline{toc}{chapter}{Resumen}
% put your text here


El secado solar del café constituye una alternativa energéticamente sostenible para reducir el contenido de humedad del grano; sin embargo, su desempeño se encuentra condicionado por la variabilidad de la radiación solar, la temperatura ambiente, la humedad relativa, el viento y las características térmicas y geométricas del secador. Estas perturbaciones pueden producir distribuciones no uniformes de temperatura y humedad, prolongar el tiempo de secado y generar diferencias en las condiciones experimentadas por el producto. En respuesta a esta problemática, la presente investigación propone el desarrollo de un marco integrado para el modelado, caracterización y control inteligente de las condiciones termo-higrométricas en secadores solares de café.\\

El sistema se concibe como una planta dinámica, abierta, no lineal y multivariable, en la cual interactúan la radiación solar, la conducción en los componentes sólidos, la convección natural o forzada, el transporte de vapor y la evaporación de humedad desde el producto. Para su estudio se plantea una metodología multiescala que combina modelos físicos de parámetros concentrados, simulaciones de dinámica de fluidos computacional, mediciones experimentales distribuidas y modelos basados en datos. El modelo CFD permite analizar los campos espaciales de temperatura, velocidad y humedad, mientras que las representaciones reducidas e híbridas se orientan a la predicción rápida, la estimación de estados y su eventual incorporación en estrategias de control.\\

La investigación considera estrategias de control clásico, control predictivo basado en modelos y control apoyado en modelos híbridos físico–data-driven. El objetivo de control no se limita al seguimiento de valores promedio de temperatura y humedad, sino que incorpora la reducción de su variabilidad espacial, las restricciones físicas de los actuadores y la influencia anticipada de las perturbaciones meteorológicas. Asimismo, se plantea la identificación de configuraciones suficientes de observación y actuación que permitan reducir la complejidad tecnológica sin comprometer la capacidad de regulación del proceso.\\

La contribución propuesta reside en la integración coherente de modelos con diferentes niveles de fidelidad dentro de una arquitectura común de observación, validación, predicción y control. Esta arquitectura distingue explícitamente entre el dominio respaldado experimentalmente, el dominio explorado mediante simulación y las condiciones de extrapolación, con el propósito de evitar que la capacidad predictiva de los modelos sea asumida fuera de las regiones sustentadas por evidencia. Se espera que el marco desarrollado contribuya al diseño de sistemas de secado solar más uniformes, adaptables y computacionalmente viables para las condiciones ambientales de Panamá.




\bigskip
 
\noindent\textit{Panamá, Agosto 2026}
\hfill A.~R.
